% Terminal note of the homily: everything the speech itself must not carry ---
% the audience, the spoken word count and the pace it implies, the relation to
% the three readings of the expansive study, the route by which each quoted
% text reaches the page, and the exact loci --- followed by the References.
% Nothing here is read aloud.

\section*{Note on Sources and Delivery}
\runninghead{Sources and Delivery}

{\footnotesize

\textbf{Audience and occasion.} An adult parish assembly, at the Mass
\latin{Dominica vigesima prima post Pentecosten}, II classis, of the 1962
\work{Missale Romanum}, editio typica, pp.~413--415, nos.~1699--1708, for the
occurrence of 18~October 2026. The oration \latin{pro Fidei propagatione},
added that day, and the commemoration of St~Luke in a low or conventual Mass
belong to other formularies and are not used.

\textbf{Length and pace.} The spoken body runs to 1,442 words, about 11.1 to
12.0 minutes at an unhurried 120 to 130 words a minute by arithmetic alone; no
one has delivered or timed it. It was read through in full, silently, for sense
and ease of speech.

\textbf{Relation to the reviewed readings.} The homily follows the second of
the full study's three readings, \textit{The King's Reckoning and the Forgiven
Debtor}: the two debts, the gift beyond the request (Chrysostom), every man a
debtor with a debtor and the forgiven debtor going out from Baptism
(Augustine, whose own join this is), the plea repeated, the pardon revoked by
its receiver (Chrysostom), Aquinas's teaching that forgiven sins return only
through ingratitude, Augustine on tongue and heart and on correction without
rancour, and his sermon on the Lord's Prayer beside the Communion's plea. From
the first reading, \textit{The War God Permits and Bounds}, it takes the
Introit's will that none resists, the Offertory's petition and permission, and
Chrysostom's wrath poured on the devil, his own join of the wronged man and the
devil (\work{Homilies on Ephesians} 22, in the moral part on Eph 6:10--13). The
third reading is left to the studies. Following the Homiletic Directory, the
speech proclaims the debt forgiven in Christ's death and Resurrection and
received in Baptism, and prepares the assembly for the Eucharist as a share in
that mystery, before it draws the consequence for the brother. These joins are
the homily's own and credited to no Father: the remission located on the Cross
and at the empty tomb; the Secret's offerings as the one sacrifice of the Cross
made present; Augustine's beggars who ask God for Christ the living Bread (said
of prayer, not of Communion) set beside the altar; the food of immortality as a
share in the Lord's death and rising; the Lord's Prayer, said before the
Communion, as the hinge between altar and brother; the wrong throat; the
Introit and Offertory as the measure of the enemy's reach; the Collect's family
as the household of forgiven servants; and the Communion's plea answered from a
sermon that speaks of the martyrs and not of this chant. The Gradual and
Alleluia are left to the studies.

\textbf{Texts quoted, and their route.} Scripture is the Douay--Rheims,
Challoner's revision: Mt 18:26, 28, 29, 32, 34 and 35; Eph 6:12; Mt 20:2 is
paraphrased. The Introit's first sentence, the Collect's first words, the
Secret, the Communion's question and the Postcommunion's ``food of
immortality'' are in the anonymous English printed by Eugene Cummiskey at
Philadelphia in 1861, a historical study aid and not an approved liturgical
translation of the 1962 Missal. The Offertory is paraphrased; the
Postcommunion's petition is given from the Latin \latin{sectemur}, ``follow
after'', which that book renders ``receive''. The Lord's Prayer is quoted as
Augustine's translator gives it. Chrysostom and Augustine are quoted in the
English of the Nicene and Post-Nicene Fathers; Aquinas is paraphrased from his
Latin.

\textbf{Exact loci.} Introit, no.~1699; Collect, no.~1700; Epistle, Eph 6:12,
no.~1701; Gospel, Mt 18:23--35, no.~1704; Offertory, no.~1705; Secret,
no.~1706; Communion, Ps 118:84 (119:84), no.~1707; Postcommunion, no.~1708.
Chrysostom, \work{Homilies on Matthew} 61, 3 (the gift beyond the request) and
4 (``straightway'', the harbour, the pardon revoked by its receiver);
\work{Homilies on Ephesians} 22. Augustine, \work{Sermon} 83 (NPNF Sermon 33),
2 (every man a debtor; God's beggars and the living Bread, paraphrased), 4 (the
Lord's Prayer), 7 (Baptism; tongue and heart) and 8 (the surgeon);
\work{Sermon} 58 (NPNF Sermon 8, on the Lord's Prayer), 8. Aquinas,
\work{Super Matthaeum} c.~18, on v.~35.

\par}

\Needspace{14\baselineskip}
\section*{References}
\runninghead{References}

{\footnotesize
\begin{itemize}\setlength{\itemsep}{0.15em}\setlength{\parskip}{0pt}
\item \work{Missale Romanum ex decreto Sacrosancti Concilii Tridentini restitutum, Summorum Pontificum cura recognitum}, editio typica (Typis Polyglottis Vaticanis, 1962), CMAA facsimile, pp.~413--415, nos.~1699--1708.
\item \work{The Roman Missal, Translated into the English Language for the Use of the Laity} (Philadelphia: Eugene Cummiskey, 1861), pp.~450--452 (XXI Sunday after Pentecost: Introit, Collect, Offertory, Secret, Communion, Postcommunion).
\item The Holy Bible, Douay--Rheims version, Challoner's revision: Mt 18:23--35; 20:2; Eph 6:12.
\item St John Chrysostom, \work{Homilies on Matthew} 61, Nicene and Post-Nicene Fathers, 1st ser., vol.~10 (CCEL text); \work{Homilies on Ephesians} 22, Nicene and Post-Nicene Fathers, 1st ser., vol.~13 (CCEL text).
\item St Augustine, \work{Sermons} 58 and 83 (Benedictine numbering; Sermons 8 and 33 in the translation), Nicene and Post-Nicene Fathers, 1st ser., vol.~6 (CCEL text).
\item St Thomas Aquinas, \work{Super Evangelium S.~Matthaei lectura} c.~18, in \work{Opera}, editio altera Veneta, vol.~3 (Venice, 1745), pp.~242--244; paraphrased, not quoted.
\end{itemize}
\par}
